\subsection{Individual differences in behavior are mirrored in population geometry}

\begin{figure}
    \centering
     \includegraphics[width=0.99\textwidth]{content/figures/ibl.pdf}

    % \caption{   \textbf{Behavioral strategy and genetic background are reflected in neural geometry.}
    % \textbf{a.} IBL dataset schematics.
    % \textbf{b.} Psychometric curves all animals (n=56) in gray and grand average in blue.
    % \textbf{c.} Matrix of neural and behavioral dissimilarities between all animals. Note that in all the analyses, independent data was used to compute behavioral and neural distances. Here shown for one data split.  \textbf{d.} Left, correlation between behavioral and neural distances of one fold. Right, correlation values of n=100 repetitions of the data split and similar analyses for shuffled data. Both done using neural data around the response ([0-0.4 s]). \textbf{e.} The correlation between neural and behavioral distances throughout the trial and aligned to stimulus or response. Red, blue and black ticks on the top mark the time-points where stimulus, response or stimulus vs response was significantly different than 0 (bootstrap t-test), respectively. \textbf{f.} Alignment between neural representations allows for 0-shot learning of psychometric curves using the 2-closest neighbors. Behavioral correlation between the first (2 closest neighbors) and last point (2 furthest neighbors) is significantly different ($t=2.2, p=0.037$). }
    \caption{\textbf{Behavioral strategy and genetic background are reflected in neural geometry.}
      \textbf{a.} IBL dataset schematics. Inset: Mean shape distance between the population responses of five recorded regions (CA1, DG, LP, PO, VISa), averaged over the $n = 17$ mice recorded in all five. Wilcoxon signed-rank across mice; $^{*}p <
  0.05$, $^{**}p < 0.01$, $^{***}p < 0.001$.
      \textbf{b.} Psychometric curves for all $n = 56$ mice. Grey, individual mice; teal, group mean.
      \textbf{c.} Pairwise distance matrices across the 56 mice. Lower triangle, neural shape distance (greys); upper triangle, behavioural distance (purples). 
      \textbf{d.} Neural against behavioural distance for all 1{,}540 mouse pairs.
      \textbf{e.} Distribution across cross-validation folds of the z-scored correlation (against shuffled distributions) between neural and behavioural distance, per region  filled circle, mean; whiskers, 95\% C.I.
      \textbf{f.} Behaviour of mice grouped by neural geometry alone. Neural distance matrix in \textit{c)} was clustered with optimal $k = 3$ (\hl{sup}) Mean $\pm$ s.e.m.
  psychometric curve of cluster 1 ($n = 24$) and cluster 2 ($n = 30$); cluster 3 ($n = 2$) is omitted, as it consisted of just two animals (\hl{sup})
      \textbf{g.} Genotype comparison. Left, dendrogram, Ward linkage with optimal leaf ordering on the
  neural distance matrix; Right, the same mice embedded by MDS.
      \textbf{i.} For each wild-type mouse, its mean neural distance to the other wild-type mice and to each mutant line. Errorbars are $\pm$ s.e.m. and $p$ values from a permutation test
  over genotype labels (10{,}000 permutations).}

    \label{fig:ibl}
\end{figure}

Individual differences in geometry need not be confined to a single region: we next asked whether they track an animal's behavior, and whether such a link is localized to particular brain regions or distributed across them.
To do this, we turned to public data released by the International Brain Lab \parencite{iblreproducibility}.
Two features of this dataset make it an ideal candidate for this analysis.
First, it consists of neural and behavioral measurements from a large cohort of animals performing the same binary choice task involving contrast discrimination (\textit{Fig.} \ref{fig:ibl}a).
Second, a large number of neurons were measured in each animal spanning the same brain regions including secondary visual areas, hippocampus, and thalamus.

We included mice with more than n = 100 neurons and n = 40 behavioral trials (i.e. n = 20 trials per cross-validation fold), yielding a dataset of $K= 56$ subjects. As an input to our shape analyses, we computed 90-dimensional tuning curves (n = 9 contrast conditions and n = 10 time bins of 40 ms after the animals response) for a total of n=24583 neurons. As done in previous analyses, we could detect differences across all recorded regions using Procrustes distances (\textit{Fig.} \ref{fig:ibl}a, inset), despite tat regions recorded in this dataset have previously been shown to encode all task variables \parencite{iblreproducibility}.
To quantify the behavioral profiles on the task, we computed the probability of answering 'right' $P(\text{right})$ for each contrast.
These curves show considerable variability across subjects (\textit{Fig.} \ref{fig:ibl}b), prompting us to ask whether there is a neural signature of this variability.

We started by computing matrices of pairwise distances across subjects quantifying neural similatiry using Procrustes \cref{eq:procrustes-distance}) and behavioral similarity using Pearson correlation (Fig. \ref{fig:ibl}c). A correlation analysis between behavioral and neural distances (\textit{Fig.} \ref{fig:ibl}c) revealed a significant relationship between these two measures (\textit{Fig.} \ref{fig:ibl}d). Repeating this analysis using independent trials to calculate neural and behavioral distances between animals (thus avoiding spurious correlations due to shared noise) confirmed that the correlation values were consistently higher than the correlations between shuffled subjects (\textit{Fig.} \ref{fig:ibl}d). This analysis confirmed that behavioral strategies are reflected in the encoding geometry. 

To pinpoint which region drove this relationship, we performed similar analyses conditioning on 5 different regions: visual area (VISa), CA1, dentate gyrus (DG) , and lateral posterior (LP) and posterior (PO) thalamic nuclei. Specifically, we examined the correlation between neural and behavioral distances, as above, but using neural activity from each region separately, using mice with more than 40 trials in the analyzed region. \hl{We found significant differences for PO ($p=0.03$, but not for the others (VISa, $p=0.17$; CA1, $p=0.28$; DG, $p=0.10$; LP, $p=40$).} This analysis suggested that despite noticeable geometric differences between regions (\textit{Fig.} \ref{fig:ibl}a, inset), the relationship between behavioral strategies and encoding geometry was distributed across most regions. \hl{missing description of fig 6f}

\subsection{Distinguishing neural representations across mice genotypes}
Individual differences can also be inherited. In a final analyses, and at the finest scale of individual variability we consider, we asked whether genetic background drives population geometry, even where single-neuron analyses report no difference (\cite{noel2025}). We analyzed a cohort in which the same visual discrimination task was performed by
wild-type mice and by three mouse lines carrying mutations associated with autism
spectrum disorder \hl{(\textit{C57BL/6}, $n=15$; \textit{Fmr1}, $n$ = 21; \textit{Cntnap2}, $n$ = 20 and \textit{Shank3}, $n$ = 19)} released by
\cite{noel2025}
That study characterized these genotypes through the lens of prior encoding. Stimulus processing, by contrast, was reported to be unaffected: psychometric performance did not differ across genotypes, and inspection of example rasters and peri-stimulus time histograms showed the same standard response features in mouse every line.
That comparison, however, was made at the level of single neurons. Population geometry is a different level of description \parencite{Kriegeskorte2021}, and whether it also is indistinguishable across genotypes has not been asked. Moreover, an important question remains open: the original analyses pools units within a genotype, so no individual animal is ever compared to another individual animal, arguably the natural unit for asking whether the three lines differ from controls.

For each animal, we assembled a matrix of stimulus-evoked encoding kernels across neurons, restricted to the three lowest-contrast conditions and 106 time bins (as in \cite{noel2025}), and retained animals with at least n = 100 units after outlier removal, \hl{leaving $n=9$ control animals and $n=9$, $n=6$, and $n=9$ for \textit{Fmr1}, \textit{Cntnap2}, and \textit{Shank3}, respectively}. We then computed a $K \times K$ matrix of pairwise shape distances across all animals using Procrustes and pooling genotypes, so that each animal occupies a single position in a common space. Hierarchical clustering of this matrix suggested that wild-type animals are clustered together (\textit{Fig.} \ref{fig:ibl}g), and the same structure was visible in a 2D embedding obtained by MDS followed by PCA (\textit{Fig.} \ref{fig:ibl}h). To quantify this, we compared for each wild-type animal its mean distance to the remaining wild-type animals against its mean distance to each mutant line (\textit{Fig.} \ref{fig:ibl}c). Indeed, distances between control animals were significantly smaller than between controls and mutants (\hl{$p= XXXX$}; see \textit{Fig.} \ref{fig:ibl}i for pairwise tests). Moreover, the three mutant lines were interchangeable when comparing each genotype against other mutants (\textit{Shank3} vs other mutants, $p=0.21$; \textit{Fmr1} other mutants, $p=0.51$; \textit{Cntnap2} other mutants, $p=0.97$; permutation test). This last observation is in line with the picture of convergent phenotypes suggested by \cite{noel2025}: distinct genetic perturbations may converge on a common computational anomaly, reflected here in similar population geometry.

Thus population geometry distinguishes genotypes that appear identical at the single-neuron level, mirroring what we found across regions, where geometry recovered an anatomical hierarchy that single-variable decodability had obscured. In both cases, structure invisible single-neuron analysis becomes visible once neural systems are compared as geometric objects.
%Finally, we show once again (see also Figs. \ref{fig:synth} and \ref{fig:hdcells}) how aligning neural representations allows for zero-shot learning of psychometric curves using the two closest neighbors. Specifically, for every subject in the dataset we computed the behavioral correlation between this subject and the two nearest neighbor subjects in the neural shape space (closest, in \textit{Fig.} \ref{fig:ibl}f). We then repeated this analysis, but now predicting each subject behavior using progressively distant subjects in neural shape space. We found that the behavioral correlation between the first (two closest neighbors) and the last (two furthest neighbors) was significantly different ($t=2.2, p=0.037$). Furthermore, we found a negative correlation between distance in neural shape space and behavioral distance (\textit{Fig.} \ref{fig:ibl}f), showing that the behavioral correlation decreases smoothly for increasingly distant animals. All together, these analyses show how different behavioral strategies are reflected in different neural geometries, as quantified by shape metrics.

